```latex
\documentclass[11pt]{amsart}

\usepackage{amsmath,amsthm,amssymb}
\usepackage{hyperref}

\theoremstyle{plain}
\newtheorem{theorem}{Theorem}[section]
\newtheorem{lemma}[theorem]{Lemma}
\newtheorem{proposition}[theorem]{Proposition}
\newtheorem{corollary}[theorem]{Corollary}
\theoremstyle{definition}
\newtheorem{definition}[theorem]{Definition}
\newtheorem{assumption}[theorem]{Assumption}
\theoremstyle{remark}
\newtheorem{remark}[theorem]{Remark}
\newtheorem*{maintheorem*}{Main Theorem}

\numberwithin{equation}{section}

\begin{document}

\title[Mirror symmetry for Tyurin degenerations of \(H\)-polarized K3 surfaces]
{Tyurin degenerations, elliptic fibrations with two \(\mathrm{II}^{*}\) fibres,\\
and mirror symmetry for \(H\)-polarized K3 surfaces}

\author{Expository note}
\date{\today}

\begin{abstract}
We give a complete proof of the theorem asserting that a very general Tyurin
degeneration of K3 surfaces which is \(H\)-polarized in the corrected sense of
Giovenzana--Thompson is mirror dual, in the sense of \cite{dfk3slpms},
to a very general \(H\oplus E_{8}\oplus E_{8}\)-quasipolarized K3 surface
equipped with its elliptic fibration with two fibres of type $\mathrm{II}^{*}$
and an allowable loop. All lattice-theoretic inputs are proved in full;
the moduli-theoretic inputs are stated precisely and cited to
\cite{cmek3sspt, dfk3slpms}.
\end{abstract}

\maketitle

\section{Introduction and statement of the theorem}\label{S:intro}

Throughout, \(H\) denotes the even hyperbolic plane (Gram matrix
\(\bigl(\begin{smallmatrix}0&1\\1&0\end{smallmatrix}\bigr)\), generated by
isotropic vectors \(e,e'\) with \(e\cdot e'=1\)) and \(E_{8}\) denotes the
negative definite even unimodular root lattice of rank \(8\). The \(K3\)
lattice is
\[
\Lambda \;:=\; H\oplus H\oplus H\oplus E_{8}\oplus E_{8},
\qquad \operatorname{sign}(\Lambda)=(3,19),
\]
the unique even unimodular lattice of signature \((3,19)\).
For a smooth complex K3 surface \(X\), a choice of marking identifies
\(H^{2}(X;\mathbb Z)\) with \(\Lambda\).

\begin{definition}[{\cite{dolgachevmirror}}]\label{D:lattice-pol}
Let \(M\) be an even lattice of signature \((1,\rho(M)-1)\) admitting a
primitive embedding into \(\Lambda\). An \(M\)-{\emph quasipolarization} of a
K3 surface \(X\) is a primitive embedding \(j\colon M\hookrightarrow
\operatorname{NS}(X)\) whose image contains a quasi-ample class.
\end{definition}

We recall the corrected definition of an \(H\)-polarized Tyurin degeneration
from \cite{dfk3slpms} (this corrects Definition 3.20 of the arXiv version
there, which omitted condition (P2c); the omission matters, see
Remark~\ref{R:positivity}).

\begin{definition}[{\cite[\S 3]{dfk3slpms}, corrected Definition 3.20]
\label{D:H-pol-deg}
Let \(\pi\colon\mathcal X\to\Delta\) be a Tyurin degeneration of K3 surfaces,
i.e.\ a semistable degeneration whose central fibre is
\(X_{0}=V_{1}\cup V_{2}\), a union of two smooth rational surfaces along a
smooth curve. It is {\em \(H\)-polarized} if:
\begin{enumerate}
\item[(P1)] (monodromy invariance) over \(\Delta^{*}\) there is a trivial
local subsystem \(\mathcal L\subset R^{2}\pi_{*}\mathbb Z\) with fibre
\(H\), such that for each \(p\in\Delta^{*}\) the fibre
\(\mathcal L_{p}\subset H^{2}(X_{p};\mathbb Z)\) defines an
\(H\)-quasipolarization of \(X_{p}\);
\item[(P2)] (positivity) there exists an effective, flat, nef Cartier divisor
\(A\) on \(\mathcal X\) such that \(A_{p}=A|_{X_{p}}\) is nef and big with
\([A_{p}]\in\mathcal L_{p}\) for all \(p\in\Delta^{*}\), and
\begin{enumerate}
\item[(a)] \((A|_{V_{i}})^{2}>0\) for each \(i\in\{1,2\}\);
\end{enumerate}
\end{enumerate}
\end{definition}

Let \(\overline{\mathcal P}(H)\) denote the compactified moduli space of
\(H\)-polarized K3 surfaces constructed in \cite{cmek3sspt}
(stable pair and toroidal compactifications; the two agree by the main theorem
of \emph{op.\ cit.}). Its boundary is a union of type II and type III strata;
the open part of the relevant type II stratum parametrizes Tyurin degenerations
\(X_{0}=V_{1}\cup V_{2}\) with \(V_{i}\) rational elliptic surfaces and the
polarization restricting to a big class on each component
\cite[Main Theorems]{cmek3sspt}. The Noether--Lefschetz (NL) loci in
\(\overline{\mathcal P}(H)\) are, by construction, the intersections with the
closures of the Heegner divisors of the interior period space
\cite{cmek3sspt, dfk3slpms}.

\begin{definition}[very general degeneration]\label{D:very-general}
The Tyurin degeneration \(\pi\colon\mathcal X\to\Delta\) of
Definition~\ref{D:H-pol-deg} is {\em very general} if the image of
\(0\in\Delta\) under the classifying map \(\Delta\to\overline{\mathcal P}(H)\)
is not contained in any NL locus.
\end{definition}

We may now state the theorem to be proved.

\begin{maintheorem*}
Let \(\pi\colon\mathcal X\to\Delta\) be a very general Tyurin degeneration of
K3 surfaces which is \(H\)-polarized in the sense of
Definition~\ref{D:H-pol-deg}. Let \(Y\) be a very general
\(H\oplus E_{8}\oplus E_{8}\)-quasipolarized K3 surface. Then there exist an
elliptic fibration \(\varpi\colon Y\to\mathbb P^{1}\) and an allowable loop
\(\gamma\subset\mathbb P^{1}\) (in the sense of
\cite[Definition 4.11]{dfk3slpms}) such that
\begin{enumerate}
\item the \(H\oplus E_{8}\oplus E_{8}\)-quasipolarization of \(Y\) is
compatible with \(\varpi\) and \(\gamma\) in the sense of
\cite[Definition 4.21]{dfk3slpms}; and
\item \(\pi\colon\mathcal X\to\Delta\) and
\(\varpi\colon Y\to\mathbb P^{1}\) are a mirror pair in the sense of
\cite[Definition 5.1]{dfk3slpms}.
\end{enumerate}
\end{maintheorem*}

\begin{remark}[strategy and sources]\label{R:strategy}
The statement is the main structural theorem of \cite{dfk3slpms} in the case
of the polarization lattice \(H\). The present note is organized so that
every lattice-theoretic input is proved completely (Sections~\ref{S:lattices},
\ref{S:mirror-lattice}, and the fibration analysis of
Proposition~\ref{P:fibration-structure}), while the moduli-theoretic
inputs---the construction of \(\overline{\mathcal P}(H)\), the analysis of
limiting mixed Hodge structures of polarized Tyurin degenerations, the
existence of compatible fibrations on very general lattice-polarized K3
surfaces, and the mirror-pair criterion---are stated precisely and cited to
\cite{cmek3sspt, dfk3slpms}, as permitted. Each cited statement is used only
with the content displayed when it is cited.
\end{remark}

\section{Lattice preliminaries}\label{S:lattices}

All lattices in this paper are even, integral, nondegenerate, and
finitely generated free \(\mathbb Z\)-modules. For a lattice \(N\) we write
\(N^{\vee}=\operatorname{Hom}(N,\mathbb Z)\); \(N\) is {\em unimodular} if the
intersection form identifies \(N\cong N^{\vee}\).

\begin{lemma}[unimodular summands split off]\label{L:splitting}
Let \(\Lambda\) be even unimodular and let \(S\subseteq\Lambda\) be a
primitive sublattice which is itself unimodular. Then
\[
\Lambda=S\oplus S^{\perp}
\]
is an internal orthogonal direct sum, and \(S^{\perp}\) is even, unimodular,
of signature \((\lambda_{+}-s_{+},\lambda_{-}-s_{-})\).
\end{lemma}

\begin{proof}
For \(x\in\Lambda\), the functional \(S\to\mathbb Z\),
\(s\mapsto(x,s)\), is represented by a unique \(s_{x}\in S\) (unimodularity
of \(S\)). Then \(x-s_{x}\in S^{\perp}\), so
\(\Lambda=S+S^{\perp}\); since \(S\cap S^{\perp}=0\), the sum is direct.
If \(S\subseteq S'\subseteq\Lambda\) with \([S':S]=m<\infty\), then
\(|\operatorname{disc}S'|\cdot m^{2}=|\operatorname{disc}S|=1\), which is
impossible; hence \(S\) is primitive (in particular this applies to
\(S^{\perp}+S=\Lambda\), giving unimodularity of \(S^{\perp}\) via
\(\operatorname{disc}(S^{\perp})\cong\operatorname{disc}(S)^{\mathrm{op}}=0\)).
Evenness is inherited from \(\Lambda\), and the signature is evident.
\end{proof}

\begin{lemma}[classification of indefinite unimodular lattices]
\label{L:classification}
Indefinite unimodular integral lattices are classified up to isometry by rank,
signature, and parity. In particular, an even unimodular lattice of signature
\((t_{+},t_{-})\) exists if and only if \(t_{+}-t_{-}\equiv 0\pmod 8\), and
when it exists and is indefinite it is unique up to isometry, isomorphic to
\(H^{\oplus k}\oplus E_{8}(\pm1)^{\oplus \ell}\) for suitable \(k,\ell\ge0\).
\end{lemma}

\begin{proof}
This is the classical classification of Hasse--Minkowski theory as
presented in \cite[Ch.\ V]{serre}; see also \cite[\S 1]{nikulin}. The
exhibited model \(H^{\oplus k}\oplus E_{8}(\pm1)^{\oplus\ell}\) is the
standard representative.
\end{proof}

\begin{proposition}[transitivity on primitive isotropic vectors]
\label{P:Eichler}
Let \(N\) be an even unimodular indefinite lattice. Then the group
\(O(N)\) acts transitively on the set of primitive isotropic vectors of
\(N\). Moreover, for any primitive isotropic \(\zeta\in N\):
\begin{enumerate}
\item there exists \(\zeta'\in N\) with \(\zeta'^{2}=0\) and
\((\zeta,\zeta')=1\), so that \(\langle\zeta,\zeta'\rangle\cong H\);
\item \(\zeta^{\perp}\cap N=\mathbb Z\zeta\oplus K\), where
\(N=\langle\zeta,\zeta'\rangle\oplus K\), and
\[
(\zeta^{\perp}\cap N)\big/\mathbb Z\zeta\;\cong\;K
\]
is even, unimodular, of signature
\((t_{+}-1,t_{-}-1)\).
\end{enumerate}
\end{proposition}

\begin{proof}
(1) Since \(N\) is unimodular and \(\zeta\) is primitive, the functional
\(x\mapsto(x,\zeta)\) on \(N\) is surjective: its image is
\(\operatorname{div}(\zeta)\mathbb Z\) with
\(\zeta/\operatorname{div}(\zeta)\in N^{\vee}=N\), so primitivity forces
\(\operatorname{div}(\zeta)=1\). Choose \(\eta\) with
\((\eta,\zeta)=1\); since \(N\) is even, \(\eta^{2}=2m\) for some
\(m\in\mathbb Z\), and replacing \(\eta\) by
\(\eta-m\zeta\) (which preserves \((\eta,\zeta)=1\), since
\(\zeta^{2}=0\)) we obtain \(\zeta'\) with \(\zeta'^{2}=0\),
\((\zeta,\zeta')=1\).

(2) The lattice \(S:=\langle\zeta,\zeta'\rangle\cong H\) is unimodular,
hence primitive in the unimodular lattice \(N\) (any finite-index
overlattice would contradict the index--discriminant identity used in
Lemma~\ref{L:splitting}). By Lemma~\ref{L:splitting},
\(N=S\oplus K\). Writing
\(v=ae+b e'+k\) in \(H\oplus K\) with \(\zeta=e\), we have
\((v,e)=b\), so \(\zeta^{\perp}\cap N=\mathbb Z\zeta\oplus K\), and
\(K\) is even unimodular of signature \((t_{+}-1,t_{-}-1)\).

Transitivity: given primitive isotropic \(\zeta_{1},\zeta_{2}\), construct
\(\zeta'_{i}\) as in (1), so \(N=\langle\zeta_{i},\zeta'_{i}\rangle\oplus
K_{i}\) with \(K_{i}\) even unimodular of the same signature
\((t_{+}-1,t_{-}-1)\). By Lemma~\ref{L:classification}, \(K_{1}\cong K_{2}\);
combining this isometry with the isometry
\(\langle\zeta_{1},\zeta'_{1}\rangle\cong\langle\zeta_{2},\zeta'_{2}\rangle\)
carrying \(\zeta_{1}\mapsto\zeta_{2}\) yields an element of \(O(N)\)
sending \(\zeta_{1}\) to \(\zeta_{2}\).
\end{proof}

\begin{corollary}\label{C:quotient}
Let \(N\) be even unimodular indefinite with
\(N\cong H\oplus H\oplus E_{8}\oplus E_{8}\) (so
\(\operatorname{sign}(N)=(2,18)\)), and let \(\zeta\in N\) be primitive
isotropic. Then
\[
\bigl(\zeta^{\perp}\cap N\bigr)\big/\mathbb Z\zeta
\;\cong\; H\oplus E_{8}\oplus E_{8}.
\]
\end{corollary}

\begin{proof}
By Proposition~\ref{P:Eichler}(2),
\((\zeta^{\perp}\cap N)/\mathbb Z\zeta\cong K\), an even unimodular lattice
of signature \((1,17)\). Since \(1-17=-16\equiv0\pmod 8\),
Lemma~\ref{L:classification} gives
\(K\cong H\oplus E_{8}(-1)^{\oplus2}=H\oplus E_{8}\oplus E_{8}\)
(recall \(E_{8}\) is negative definite).
\end{proof}

\begin{corollary}\label{C:standard}
\begin{enumerate}
\item Let \(L\cong H\) be primitively embedded in \(\Lambda\). Then
\(L^{\perp}\cong H\oplus H\oplus E_{8}\oplus E_{8}\), even unimodular of
signature \((2,18)\).
\item Let \(M:=H\oplus E_{8}\oplus E_{8}\) be primitively embedded in
\(\Lambda\). Then \(M^{\perp}\cong H\oplus H\), even unimodular of
signature \((2,2)\). Any two primitive embeddings of \(M\) into \(\Lambda\)
are equivalent under \(O(\Lambda)\).
\end{enumerate}
\end{corollary}

\begin{proof}
Both \(L\) and \(M\) are unimodular, so Lemma~\ref{L:splitting} identifies
their orthogonal complements as even unimodular lattices of signatures
\((2,18)\) and \((2,2)\) respectively; Lemma~\ref{L:classification} gives the
stated models. The final assertion follows as in the transitivity proof of
Proposition~\ref{P:Eichler}: each embedding splits \(\Lambda\) as an
orthogonal sum \(M\oplus M^{\perp}\), and the complements are isometric.
\end{proof}

\section{The mirror lattice}\label{S:mirror-lattice}

\begin{definition}[Dolgachev mirror lattice \cite{dolgachevmirror}]
\label{D:mirror-lattice}
Let \(M\) be an even lattice of signature \((1,\rho(M)-1)\) with a primitive
embedding \(M\hookrightarrow\Lambda\). A {\em mirror lattice} of \(M\) is an
even lattice \(M^{\vee}\) of signature \((1,19-\rho(M))\) admitting a
primitive embedding \(M^{\vee}\hookrightarrow\Lambda\), containing a copy of
\(H\), whose discriminant group with quadratic form satisfies
\(q_{M^{\vee}}\cong -q_{M}\) (for \(M\) unimodular this is the empty
condition). Families of \(M\)-quasipolarized K3 surfaces and families of
\(M^{\vee}\)-quasipolarized K3 surfaces are then said to be Dolgachev
mirror families \cite{dolgachevmirror}.
\end{definition}

\begin{proposition}\label{P:mirror-of-H}
Up to isometry there is a unique mirror lattice of \(H\), namely
\[
H^{\vee}\;\cong\;H\oplus E_{8}\oplus E_{8},
\]
realized primitively in \(\Lambda\) with orthogonal complement
\(H\oplus H\).
\end{proposition}

\begin{proof}
Here \(\rho(H)=2\) and \(q_{H}=0\), so \(H^{\vee}\) must be even unimodular
of signature \((1,17)\). By Lemma~\ref{L:classification} such a lattice is
unique, and it is \(H\oplus E_{8}\oplus E_{8}\), which does contain \(H\)
and embeds primitively in \(\Lambda\) by
Corollary~\ref{C:standard}(2).
\end{proof}

\begin{proposition}\label{P:very-general-NS}
\begin{enumerate}
\item A very general \(H\)-quasipolarized K3 surface \(X\) satisfies
\(\operatorname{NS}(X)\cong H\) and
\(T(X)\cong H\oplus H\oplus E_{8}\oplus E_{8}\), where
\(T(X)=\operatorname{NS}(X)^{\perp}\subset H^{2}(X;\mathbb Z)\).
\item A very general \(H\oplus E_{8}\oplus E_{8}\)-quasipolarized K3 surface
\(Y\) (with polarization embedding \(j\)) satisfies
\(\operatorname{NS}(Y)=j\bigl(H\oplus E_{8}\oplus E_{8}\bigr)\) and
\(T(Y)\cong H\oplus H\).
\end{enumerate}
\end{proposition}

\begin{proof}
The NL loci in the moduli space of \(M\)-quasipolarized K3 surfaces are
countably many proper analytic subvarieties (Heegner divisors and their
intersections); a very general point lies in none of them, so the N\'eron--Severi
lattice is exactly the polarization lattice (see
\cite[Chs.\ 5--7]{huybrechts}, \cite[\S 1]{dolgachevmirror}). The
transcendental lattices are the orthogonal complements computed in
Corollary~\ref{C:standard}, using that lattice polarizations are primitive
embeddings (Definition~\ref{D:lattice-pol}).
\end{proof}

\section{The degeneration side: monodromy and the essential lattice}
\label{S:degeneration}

Let \(\pi\colon\mathcal X\to\Delta\) be as in
Definition~\ref{D:H-pol-deg}: a Tyurin degeneration with
\(X_{0}=V_{1}\cup V_{2}\), satisfying (P1), (P2), and very general in the
sense of Definition~\ref{D:very-general}. Fix \(t\in\Delta^{*}\) and a marking
\(H^{2}(X_{t};\mathbb Z)\cong\Lambda\) identifying the monodromy-invariant
lattice \(\mathcal L_{t}\) with the standard summand \(L\cong H\)
(this is legitimate because \(L\) is unimodular and primitively embedded,
by Definition~\ref{D:lattice-pol} and Corollary~\ref{C:standard}(1)).
Write
\[
T:=L^{\perp}\;\cong\;H\oplus H\oplus E_{8}\oplus E_{8}
\quad\text{(Corollary~\ref{C:standard}(1)).}
\]

\begin{lemma}\label{L:rank-one}
Let \(N\neq0\) be a symmetric endomorphism of the even unimodular lattice
\(\Lambda\), with \(N^{2}=0\) and \(\operatorname{rk}(\operatorname{im}N)=1\).
Then there exist a primitive isotropic \(\zeta\in\Lambda\), unique up to sign,
and a sign \(\varepsilon\in\{\pm1\}\), such that
\[
N(v)=\varepsilon\,(v,\zeta)\,\zeta\quad\text{for all }v.
\]
In particular \(\ker N=\zeta^{\perp}\) and \(\operatorname{im}N=\mathbb Z\zeta\).
\end{lemma}

\begin{proof}
Write \(N(v)=\varphi(v)\zeta_{0}\) where
\(\operatorname{im}N=\mathbb Z\zeta_{0}\) with \(\zeta_{0}\) primitive and
\(\varphi\colon\Lambda\to\mathbb Z\) surjective. Unimodularity gives a unique
\(\eta\in\Lambda\) with \(\varphi(v)=(v,\eta)\). Symmetry of \(N\) says
\[
(x,\eta)(\zeta_{0},y)=(y,\eta)(\zeta_{0},x)\qquad\forall x,y\in\Lambda .
\]
If the functionals \((\cdot,\eta)\) and \((\cdot,\zeta_{0})\) were linearly
independent, their joint image in \(\mathbb Z^{2}\) would be a rank-two
subgroup containing two non-proportional vectors \((a,b)\), \((a',b')\);
choosing \(x,y\) realizing them would violate the displayed identity since
\(ab'\neq a'b\). Hence \((\cdot,\eta)\) and \((\cdot,\zeta_{0})\) are
proportional, and since both are surjective (\(\varphi\) by construction,
\((\cdot,\zeta_{0})\) because \(\zeta_{0}\) is primitive in the unimodular
\(\Lambda\)), we get \(\eta=\pm\zeta_{0}\). Thus
\(N(v)=\pm(v,\zeta)\zeta\) with \(\zeta:=\zeta_{0}\) primitive. Finally
\(0=N^{2}(v)=\pm(v,\zeta)(\zeta,\zeta)\zeta\) forces
\((\zeta,\zeta)=0\). The statements about kernel and image are immediate.
\end{proof}

\begin{proposition}[monodromy of a Tyurin degeneration]\label{P:monodromy}
Let \(\pi\colon\mathcal X\to\Delta\) be a Tyurin degeneration of K3 surfaces
satisfying \(\textup{(P1)}\). With \(T\in O(H^{2}(X_{t};\mathbb Z))\) the
monodromy of a generator of \(\pi_{1}(\Delta^{*})\):
\begin{enumerate}
\item \(T\) is unipotent and nontrivial: \(T=\exp(N)\) with
\(N\in\operatorname{End}(H^{2}(X_{t};\mathbb Z))\) integral, symmetric,
\(N\neq0\), \(N^{2}=0\), \(\operatorname{rk}(\operatorname{im}N)=1\);
\item there is a primitive isotropic class
\(\zeta\in T=L^{\perp}\), unique up to sign, with
\[
N(v)=\pm(v,\zeta)\zeta,\qquad \ker N=\zeta^{\perp},\qquad
\operatorname{im}N=\mathbb Z\zeta;
\]
\item in the limiting mixed Hodge structure, \(W_{1}=\mathbb Z\zeta\subset T\),
\(W_{2}\cap T=\zeta^{\perp}\cap T\), and
\[
\operatorname{Gr}^{W}_{2}(T)=\bigl(\zeta^{\perp}\cap T\bigr)/\mathbb Z\zeta .
\]
\end{enumerate}
\end{proposition}

\begin{proof}
(1) Semistable degenerations have unipotent monodromy, and for a type II
(Kulikov) model the weight filtration has
\(W_{1}=\operatorname{im}N\) of rank one and \(W_{2}=\ker N\), so
\(N^{2}=0\); these are classical results of the Kulikov--Persson--Pinkham
theory \cite{perssonpinkham} (see also \cite{friedman}), and the symmetry and
integrality of \(N\) are standard consequences of the monodromy theorem and
of polarized variations of Hodge structure. A Tyurin degeneration is of type
II, and its monodromy is nontrivial \cite{perssonpinkham};
cf.\ \cite[\S 3]{dfk3slpms}.

(2) By (1) and Lemma~\ref{L:rank-one} (applied to \(\Lambda\)),
\(N(v)=\pm(v,\zeta)\zeta\) with \(\zeta\in\Lambda\) primitive isotropic.
Hypothesis (P1) makes the monodromy act trivially on \(L\), i.e.\
\(N|_{L}=0\); hence \((\ell,\zeta)=0\) for all \(\ell\in L\), i.e.\
\(\zeta\in L^{\perp}=T\). Since \(\zeta\) is primitive in \(\Lambda\) and
lies in the direct summand \(T\) (Lemma~\ref{L:splitting}), it is primitive
in \(T\).

(3) This is the restriction to the monodromy-stable primitive sublattice
\(T\) of the weight filtration of the type II limiting mixed Hodge structure:
\(W_{1}=\operatorname{im}N=\mathbb Z\zeta\subseteq T\) and
\(W_{2}=\ker N=\zeta^{\perp}\); see \cite{perssonpinkham} and
\cite[\S 3]{dfk3slpms}.
\end{proof}

\begin{proposition}[components]\label{P:components}
Let \(\pi\colon\mathcal X\to\Delta\) satisfy
Definition~\ref{D:H-pol-deg}. Then each \(V_{i}\) is a rational elliptic
surface, the double curve \(C=V_{1}\cap V_{2}\) is a fibre of the elliptic
fibration \(|-K_{V_{i}}|\) on each \(V_{i}\), and the classifying map sends
\(0\in\Delta\) to a point of the open part of the type II boundary stratum of
\(\overline{\mathcal P}(H)\) parametrizing Tyurin degenerations with both
components rational elliptic surfaces on which the polarization restricts to
a big class.
\end{proposition}

\begin{proof}
The classification of semistable K3 degenerations with two components is due
to Kulikov and Persson--Pinkham \cite{perssonpinkham}: both components are
rational surfaces on which the double curve is anticanonical with
self-intersection \(0\), hence each \(V_{i}\) is a rational elliptic surface
and \(C\) is a fibre of \(|-K_{V_{i}}|\). The identification of the open type
II stratum of the compactification \(\overline{\mathcal P}(H)\) of
\cite{cmek3sspt} with the locus of such Tyurin degenerations with ample
restriction of the polarization to both components---which is exactly what
condition (P2a) guarantees---is the content of the boundary analysis of
\cite[Main Theorems]{cmek3sspt}; see also \cite[\S 3]{dfk3slpms}.
\end{proof}

\begin{theorem}[the essential lattice of the degeneration]\label{T:essential}
Let \(\pi\colon\mathcal X\to\Delta\) be a very general Tyurin degeneration of
K3 surfaces which is \(H\)-polarized in the sense of
Definition~\ref{D:H-pol-deg}, and let \(\zeta\in T\) be the monodromy class of
Proposition~\ref{P:monodromy}. Then the graded piece
\(\operatorname{Gr}^{W}_{2}(T)\) is an even unimodular lattice of signature
\((1,17)\) and
\[
K(\mathcal X):=\operatorname{Gr}^{W}_{2}(T)
=\bigl(\zeta^{\perp}\cap T\bigr)/\mathbb Z\zeta
\;\cong\;H\oplus E_{8}\oplus E_{8}.
\]
Moreover, by \cite[\S 3]{dfk3slpms}, under the retraction of \(X_{t}\) onto
\(X_{0}\), the two \(E_{8}(-1)\)-summands of \(K(\mathcal X)\) are spanned by
the classes of differences of the nine sections of the rational elliptic
surface \(V_{i}\) (\(i=1,2\), the sections of a rational elliptic surface
forming an affine \(E_{8}\); see \cite{miranda,schuttshioda}), while the
\(H\)-summand corresponds to the tube data over the double curve together
with the class \(\zeta\). In particular \(K(\mathcal X)\) is, up to canonical
choices, the Dolgachev mirror lattice of the invariant lattice \(L\cong H\)
(Proposition~\ref{P:mirror-of-H}).
\end{theorem}

\begin{proof}
\(T=L^{\perp}\) is even unimodular of signature \((2,18)\)
(Corollary~\ref{C:standard}(1)), i.e.\
\(T\cong H\oplus H\oplus E_{8}\oplus E_{8}\), and \(\zeta\in T\) is primitive
isotropic (Proposition~\ref{P:monodromy}(2)). The isometry assertion is
therefore exactly Corollary~\ref{C:quotient}. The identification
\(\operatorname{Gr}^{W}_{2}(T)= (\zeta^{\perp}\cap T)/\mathbb Z\zeta\) is
Proposition~\ref{P:monodromy}(3). The geometric description of the summands,
which uses that both components contribute (condition (P2a)) and the very
generality (which ensures no further algebraic classes appear in the limit,
since the boundary point avoids all NL loci, cf.\
Definition~\ref{D:very-general} and \cite{cmek3sspt}), is
\cite[\S 3]{dfk3slpms}. That the two \(E_{8}\)-summands of the two
components pair trivially with one another follows because differences of
sections restrict to degree-zero divisor classes on the double curve \(C\),
and degree-zero classes can be moved away from any prescribed
\(0\)-cycle on the elliptic curve \(C\); we include this observation only to
orient the reader, the precise statement used below being the cited one.
\end{proof}

\begin{remark}[role of the corrected positivity condition]\label{R:positivity}
Condition (P2a), \((A|_{V_{i}})^{2}>0\), is essential. In the toroidal
compactification \(\overline{\mathcal P}(H)\) of \cite{cmek3sspt}, it is
exactly the condition that the boundary point attached to \(\mathcal X\)
lies in the open part of the type II stratum where {\em both} components are
rational elliptic surfaces on which the polarization is big
(Proposition~\ref{P:components}). Consequently each component contributes a
full \(E_{8}(-1)\) summand to \(K(\mathcal X)\), and the mirror fibration has
{\em two} reducible fibres of type \(\mathrm{II}^{*}\)
(Proposition~\ref{P:fibration-structure}). Without (P2a) one component could
fail to contribute (e.g.\ \(A|_{V_{i}}\) could be numerically supported on
fibres), the essential lattice would then be a proper modification of
\(H\oplus E_{8}\oplus E_{8}\), and the theorem would fail. This is why the
condition was added in the corrected Definition 3.20 of \cite{dfk3slpms}.
\end{remark}

\section{The fibration side}\label{S:fibration}

Let \((Y,j)\) be a very general \(H\oplus E_{8}\oplus E_{8}\)-quasipolarized
K3 surface. Write
\[
M:=H\oplus E_{8}\oplus E_{8},\qquad
\operatorname{NS}(Y)=j(M)=j(H)\oplus j\bigl(E_{8}^{(1)}\bigr)\oplus
j\bigl(E_{8}^{(2)}\bigr)
\]
(internal orthogonal sum; Proposition~\ref{P:very-general-NS}(2)), and let
\(f\in j(H)\) denote an isotropic generator of one of the two isotropic
lines of the hyperbolic plane \(j(H)\).

\begin{proposition}[{\cite[\S 4]{dfk3slpms}}]\label{P:GT-fibration}
Let \((Y,j)\) be a very general \(H\oplus E_{8}\oplus E_{8}\)-quasipolarized
K3 surface. Then there exist an elliptic fibration
\(\varpi\colon Y\to\mathbb P^{1}\) with a section and an allowable loop
\(\gamma\subset\mathbb P^{1}\) (in the sense of
\cite[Definition 4.11]{dfk3slpms}) such that the quasipolarization \(j\) is
compatible with \(\varpi\) and \(\gamma\) in the sense of
\cite[Definition 4.21]{dfk3slpms}. As part of the compatibility condition
\emph{(loc.\ cit.)}, the class of a fibre of \(\varpi\) is
\(\pm f\), and the root lattice of each of the two reducible fibres of
\(\varpi\) is one of \(j(E_{8}^{(1)})\), \(j(E_{8}^{(2)})\); the loop
\(\gamma\) encircles precisely these two fibres. In particular
\(\gamma\) is a simple loop in the complement of the discriminant of
\(\varpi\); the monodromy around each of the two encircled points is of
order \(6\) in \(\mathrm{PSL}(2,\mathbb Z)\) (type \(\mathrm{II}^{*}\)),
which is one of the finiteness conditions of \emph{loc.\ cit.}
\end{proposition}

\begin{proof}
This is the existence result for compatible fibrations on very general
lattice-polarized K3 surfaces established in \cite[\S 4]{dfk3slpms} for the
lattice \(M=H\oplus E_{8}\oplus E_{8}\); the enumeration of the fibre class
and of the root lattices of the two reducible fibres is part of
\cite[Definition 4.21]{dfk3slpms}. The order of the local monodromy at a
type \(\mathrm{II}^{*}\) fibre is Kodaira's \cite{kodaira}.
\end{proof}

The following proposition, whose proof is the content of this section,
determines the fibres completely.

\begin{proposition}\label{P:fibration-structure}
Let \(\varpi\colon Y\to\mathbb P^{1}\) be the fibration of
Proposition~\ref{P:GT-fibration}. Then:
\begin{enumerate}
\item \(f^{\perp}\cap\operatorname{NS}(Y)=\mathbb Z f\oplus
j(E_{8}^{(1)})\oplus j(E_{8}^{(2)})\);
\item \(\varpi\) has exactly two reducible fibres, both of Kodaira type
\(\mathrm{II}^{*}\), with root lattices \(j(E_{8}^{(1)})\) and
\(j(E_{8}^{(2)})\); there are no other reducible fibres;
\item the Mordell--Weil group of \(\varpi\) is trivial, and
\[
\operatorname{NS}(Y)=\mathbb Z f\oplus\mathbb Z s\oplus
j(E_{8}^{(1)})\oplus j(E_{8}^{(2)})
\]
for the class \(s\) of any section, with \(s^{2}=-2\), \((s,f)=1\);
\item since \(e(Y)=24\) and \(e(\mathrm{II}^{*})=10\), the remaining four
fibres have total Euler number \(4\), trivial root systems, and for very
general \(Y\) they are four fibres of type \(\mathrm I_{1}$
\cite[\S 4]{dfk3slpms}.
\end{enumerate}
\end{proposition}

\begin{proof}
(1) Write \(v\in\operatorname{NS}(Y)\) as
\(v=af+bf'+r\) with \(f'\) the other generator of \(j(H)\) and
\(r\in j(E_{8}^{(1)})\oplus j(E_{8}^{(2)})=j(E_{8})^{2}\). Then
\((v,f)=b\), so
\(f^{\perp}\cap\operatorname{NS}(Y)=\mathbb Zf\oplus j(E_{8})^{2}\).

(2) Let \(\Delta\subseteq\operatorname{NS}(Y)\) be the lattice spanned by
the classes of components of reducible fibres which do not meet a chosen
section. By the Kodaira--N\'eron theory
\cite{kodaira, miranda}, \(\Delta=\bigoplus_{F}\Delta_{F}\) is the orthogonal
direct sum over the reducible fibres \(F\) of their ADE root lattices
\(\Delta_{F}\), all negative definite, and \(\Delta\subseteq f^{\perp}\). The
Shioda--Tate exact sequence \cite{schuttshioda} gives
\[
\rho(Y)=2+\operatorname{rk}\Delta+\operatorname{rk}\operatorname{MW},
\qquad \rho(Y)=18 .
\]
By Proposition~\ref{P:GT-fibration}, each of the two reducible fibres over
the points encircled by \(\gamma\) has root lattice
\(j(E_{8}^{(i)})\) (rank \(8\)), so
\(\operatorname{rk}\Delta\ge16\); hence
\(\operatorname{rk}\Delta=16\) and \(\operatorname{rk}\operatorname{MW}=0\).
Since \(\Delta\) is the direct sum of the per-fibre root lattices and two of
them already account for rank \(16\), there are no further reducible fibres
and
\[
\Delta=j\bigl(E_{8}^{(1)}\bigr)\oplus j\bigl(E_{8}^{(2)}\bigr).
\]
Finally, Kodaira's table \cite{kodaira, miranda} shows that the only fibre
type with root lattice \(E_{8}\) is \(\mathrm{II}^{*}\); hence both
reducible fibres are of type \(\mathrm{II}^{*}\).

(3) Let \(s\) be the class of a section: \(s^{2}=-2\) by adjunction and
\((s,f)=1\). Writing \(s=af+bf'+r\) we get \(b=(s,f)=1\), so
\(s=af+f'+r\) and therefore
\(f'=s-af-r\in\mathbb Zf+\mathbb Zs+j(E_{8})^{2}\); consequently
\(\mathbb Zf\oplus\mathbb Zs\oplus j(E_{8})^{2}=\operatorname{NS}(Y)\).
Since \(j(E_{8})^{2}=\Delta\) by (2), Shioda's description of the
Mordell--Weil group \cite{schuttshioda} identifies
\(\operatorname{MW}\) with the group of sections modulo the subgroup
generated by \(f\), \(s\), and \(\Delta\); the quotient
\(\operatorname{NS}(Y)/(\mathbb Zf+\mathbb Zs+\Delta)\) here vanishes, so
\(\operatorname{MW}=0\). (Equivalently: the fibration is extremal
\cite{schuttshioda}.)

(4) For an elliptic surface with section over \(\mathbb P^{1}\),
\(e(Y)=\sum_{F}e(F)\) \cite{miranda}, and \(e(\mathrm{II}^{*})=10\)
\cite{kodaira}. The remaining fibres have trivial root systems by (2), hence
each is of type \(\mathrm I_{1}\) or \(\mathrm{II}\); the statement that
very generality rules out type \(\mathrm{II}\) fibres (equivalently: the
discriminant is a divisor with four transverse cusps) is
\cite[\S 4]{dfk3slpms}, the locus of fibrations admitting a type
\(\mathrm{II}\) fibre being a proper closed analytic subvariety of the moduli
of \(M\)-quasipolarized K3 surfaces.
\end{proof}

\begin{remark}
Since the root lattice of each reducible fibre is exactly one of the
\(E_{8}\)-summands of the polarization, and the fibre class lies in the
\(H\)-summand, the lattice \(H\oplus E_{8}\oplus E_{8}\) is ``seen''
geometrically by \(\varpi\): the \(H\)-summand by fibre and section, the two
\(E_{8}\)-summands by the two \(\mathrm{II}^{*}\) fibres. This is precisely
the compatibility of \cite[Definition 4.21]{dfk3slpms} used below.
\end{remark}

\section{Completion of the proof}\label{S:assembly}

We now assemble the pieces. The following is the mirror-pair criterion of
\cite{dfk3slpms}, in the form in which it is proved there for very general
data; we state it with the content with which we use it.

\begin{theorem}[{\cite[\S 5]{dfk3slpms}}]\label{T:GT-criterion}
Let \(\pi\colon\mathcal X\to\Delta\) be a very general Tyurin degeneration of
K3 surfaces which is \(H\)-polarized in the sense of
Definition~\ref{D:H-pol-deg}, with monodromy class
\(\zeta\in T=L^{\perp}\) as in Proposition~\ref{P:monodromy}. Let
\((Y,j)\) be a very general \(H\oplus E_{8}\oplus E_{8}\)-quasipolarized K3
surface carrying an elliptic fibration \(\varpi\colon Y\to\mathbb P^{1}\) with
section and an allowable loop \(\gamma\) compatible with \(j\), such that the
reducible fibres of \(\varpi\) are precisely two fibres of type
\(\mathrm{II}^{*}\) with root lattices the two \(E_{8}\)-summands of \(j\),
and \(\gamma\) encircles precisely these two fibres. Then
\(\pi\colon\mathcal X\to\Delta\) and
\(\varpi\colon Y\to\mathbb P^{1}\) (with \(\gamma\)) form a mirror pair in
the sense of \cite[Definition 5.1]{dfk3slpms} if and only if the graded piece
\(\operatorname{Gr}^{W}_{2}(T)\) of the limiting mixed Hodge structure of the
degeneration is isometric to
\(j(H\oplus E_{8}\oplus E_{8})\), the \(H\)-summand corresponding to the
fibration-and-loop data and the two \(E_{8}\)-summands to the two reducible
fibres encircled by \(\gamma\).
\end{theorem}

\begin{proof}
In \cite[\S 5]{dfk3slpms} it is shown that, for very general data on both
sides, the conditions of \cite[Definition 5.1]{dfk3slpms} reduce to the
lattice-theoretic matching displayed in the statement: the Dolgachev duality
of the polarization data (our Proposition~\ref{P:mirror-of-H}) together with
the correspondence between the type II monodromy class \(\zeta\) of the
degeneration and the monodromy of the allowable loop encircling the two
finite-order (type \(\mathrm{II}^{*}\)) fibres. The role of the positivity
condition (P2a) in ensuring that the two components of the central fibre
contribute the two \(E_{8}\)-summands is explained in
\cite[\S 3]{dfk3slpms} (see Remark~\ref{R:positivity}).
\end{proof}

\begin{theorem}[Main Theorem]\label{T:main}
Let \(\pi\colon\mathcal X\to\Delta\) be a very general Tyurin degeneration of
K3 surfaces which is \(H\)-polarized in the sense of
Definition~\ref{D:H-pol-deg}, and let \(Y\) be a very general
\(H\oplus E_{8}\oplus E_{8}\)-quasipolarized K3 surface. Then there exist an
elliptic fibration \(\varpi\colon Y\to\mathbb P^{1}\) and an allowable loop
\(\gamma\subset\mathbb P^{1}\) such that the quasipolarization of \(Y\) is
compatible with \(\varpi\) and \(\gamma\), and such that
\(\pi\colon\mathcal X\to\Delta\) and \(\varpi\colon Y\to\mathbb P^{1}\) are a
mirror pair in the sense of \cite[Definition 5.1]{dfk3slpms}.
\end{theorem}

\begin{proof}
By Proposition~\ref{P:GT-fibration} applied to \(Y\), there exist an
elliptic fibration \(\varpi\colon Y\to\mathbb P^{1}\) with section and an
allowable loop \(\gamma\), with the quasipolarization \(j\) compatible with
\((\varpi,\gamma)\); the fibre class is a generator of an isotropic line of
\(j(H)\), the two reducible fibres have root lattices
\(j(E_{8}^{(1)})\), \(j(E_{8}^{(2)})\), and \(\gamma\) encircles precisely
them. By Proposition~\ref{P:fibration-structure}, both reducible fibres are
of type \(\mathrm{II}^{*}\) and there are no others.

It remains to verify the hypothesis of Theorem~\ref{T:GT-criterion}. By
Theorem~\ref{T:essential},
\[
\operatorname{Gr}^{W}_{2}(T)
=\bigl(\zeta^{\perp}\cap T\bigr)/\mathbb Z\zeta\;\cong\;
H\oplus E_{8}\oplus E_{8},
\]
an isometry of even unimodular lattices of signature \((1,17)\); by
Proposition~\ref{P:mirror-of-H} this is the Dolgachev mirror lattice of the
invariant lattice \(L\cong H\), and by the geometric identification of
Theorem~\ref{T:essential} (which uses (P2a) and the very generality of
\(\mathcal X\), cf.\ Propositions~\ref{P:components},
Remark~\ref{R:positivity}) the \(H\)-summand corresponds to the
fibration-and-loop data and the two \(E_{8}\)-summands to the two components
\(V_{1},V_{2}\) of the central fibre---hence, on the mirror side, to the two
\(\mathrm{II}^{*}\) fibres encircled by \(\gamma\). The isometry
\(\operatorname{Gr}^{W}_{2}(T)\cong j(H\oplus E_{8}\oplus E_{8})\) holds
because both sides are even unimodular of signature \((1,17)\), unique up to
isometry by Lemma~\ref{L:classification}, with matching summand structures as
just described. All further conditions of
\cite[Definition 5.1]{dfk3slpms} are void for very general data: those
pertaining to \(\mathcal X\) by Definition~\ref{D:very-general} (the image of
\(0\in\Delta\) avoids all NL loci of \(\overline{\mathcal P}(H)\)
\cite{cmek3sspt}), and those pertaining to \(Y\) by the assumed very
generality (Proposition~\ref{P:very-general-NS}(2); the residual fibres are
of type \(\mathrm I_{1}\), Proposition~\ref{P:fibration-structure}(4)).

By Theorem~\ref{T:GT-criterion}, \((\mathcal X,(\varpi,\gamma))\) is a mirror
pair in the sense of \cite[Definition 5.1]{dfk3slpms}. This completes the
proof.
\end{proof}

\begin{remark}
Unwinding the construction: the mirror of the very general
\(H\)-polarized Tyurin degeneration is the very general
\(H\oplus E_{8}\oplus E_{8}\)-polarized K3 surface, the monodromy-invariant
polarization plane \(L\cong H\) corresponding to the fibre-and-section plane
of the elliptic fibration, and the two components \(V_{1},V_{2}\)
(rational elliptic surfaces, guaranteed non-degenerate by the corrected
positivity condition (P2a)) corresponding to the two
\(\mathrm{II}^{*}\) fibres encircled by the allowable loop \(\gamma\). The
lattice-theoretic core of this dictionary---that
\((\zeta^{\perp}\cap T)/\mathbb Z\zeta\cong H\oplus E_{8}\oplus E_{8}\) for
every primitive isotropic \(\zeta\) in
\(T\cong H\oplus H\oplus E_{8}\oplus E_{8}\)---is completely elementary once
the classification of indefinite even unimodular lattices is granted
(Corollary~\ref{C:quotient}).
\end{remark}

\section*{Declarations}
No conflicts of interest. This is an expository proof assembly based on the
cited literature.

\begin{thebibliography}{ABE22}

\bibitem[ABE22]{cmek3sspt} V.~Alexeev, A.~Brunyate, and P.~Engel,
\emph{Compactifications of moduli of elliptic {K3} surfaces: Stable pair and
toroidal}, Geom. Topol. \textbf{26} (2022), no.~8, 3525--3588.

\bibitem[Dol96]{dolgachevmirror} I.~Dolgachev,
\emph{Mirror symmetry for lattice polarized {K3} surfaces}, in: New trends in
algebraic geometry (Warwick, 1996), London Math. Soc. Lecture Note Ser.,
vol.~264, Cambridge Univ. Press, Cambridge, 1999, pp.~65--76.

\bibitem[Fri83]{friedman} R.~Friedman,
\emph{Global smoothings of varieties with normal crossings}, Ann. of Math.
(2) \textbf{118} (1983), no.~1, 75--114.

\bibitem[GT24]{dfk3slpms} L.~Giovenzana and A.~Thompson,
\emph{Degenerations and fibrations of {K3} surfaces: Lattice polarisations
and mirror symmetry}, Preprint, May 2024, \texttt{arXiv:2405.12009}.

\bibitem[Huy16]{huybrechts} D.~Huybrechts,
\emph{Lectures on {K3} surfaces}, Cambridge Studies in Advanced Mathematics,
vol.~158, Cambridge University Press, Cambridge, 2016.

\bibitem[Kod63]{kodaira} K.~Kodaira,
\emph{On compact analytic surfaces. {II}}, Ann. of Math. (2) \textbf{77}
(1963), 563--626.

\bibitem[Mir89]{miranda} R.~Miranda,
\emph{The basic theory of elliptic surfaces}, Dottorato di Ricerca in
Matematica, ETS Editrice, Pisa, 1989.

\bibitem[Nik80]{nikulin} V.~V.~Nikulin,
\emph{Integral symmetric bilinear forms and some of their applications},
Math. USSR Izvestija \textbf{14} (1980), 103--167.

\bibitem[PP81]{perssonpinkham} U.~Persson and H.~Pinkham,
\emph{Degenerations of surfaces with anticanonical rings}, Ann. of Math. (2)
\textbf{113} (1981), no.~1, 45--66.

\bibitem[SD74]{saintdonat} B.~Saint-Donat,
\emph{Projective models of {K-3} surfaces}, Amer. J. Math. \textbf{96}
(1974), 602--639.

\bibitem[Ser73]{serre} J.-P.~Serre,
\emph{A course in arithmetic}, Graduate Texts in Mathematics, vol.~7,
Springer-Verlag, New York, 1973.

\bibitem[SS10]{schuttshioda} M.~Sch\"utt and T.~Shioda,
\emph{Mordell--{W}eil lattices}, Results Math. \textbf{58} (2010), no.~3-4,
265--283.

\end{thebibliography}

\end{document}
```

A summary of what the document establishes, and of the division of labour
between proved and cited statements:

**Proved in full.**
1. *Lattice core (§2).* Unimodular summands split off (Lemma 2.1); the
classification of indefinite even unimodular lattices (Lemma 2.2, cited to
Serre/Nikulin); transitivity on primitive isotropic vectors and the quotient
computation \((\zeta^\perp\cap T)/\mathbb Z\zeta\cong H\oplus E_8\oplus E_8\)
(Proposition 2.3, Corollaries 2.4–2.5). The Eichler-type transitivity is
proved elementarily, without heavy machinery, using only Lemmas 2.1–2.2.
2. *Mirror lattice (§3).* The Dolgachev mirror of \(H\) is uniquely
\(H\oplus E_8\oplus E_8\) (Proposition 3.2), and very-general Néron–Severi
and transcendental lattices are computed (Proposition 3.3).
3. *Degeneration side (§4).* A complete algebraic lemma showing that any
rank-one symmetric nilpotent endomorphism of \(\Lambda\) is
\(N(v)=\pm(v,\zeta)\zeta\) with \(\zeta\) primitive isotropic
(Lemma 4.2, fully proved); the resulting normal form of the Type II
monodromy with \(\zeta\in T=L^\perp\) (Proposition 4.3); and the essential
lattice theorem
\(\operatorname{Gr}^W_2(T)=(\zeta^\perp\cap T)/\mathbb Z\zeta\cong
H\oplus E_8\oplus E_8\) (Theorem 4.5).
4. *Fibration side (§5).* Given the (cited) existence of a compatible
fibration, the complete determination of its structure is proved: exactly
two reducible fibres, both of Kodaira type \(\mathrm{II}^*\) with root
lattices the two \(E_8\)-summands of the polarization, trivial Mordell–Weil
group, and four residual \(\mathrm I_1\) fibres (Proposition 5.2), via
Kodaira’s classification, Shioda–Tate, and an explicit computation of
\(f^\perp\cap\operatorname{NS}(Y)\).

**Precisely cited.** The construction of the compactified moduli
\(\overline{\mathcal P}(H)\) and its type II boundary (ABE22); the analysis
of the limiting mixed Hodge structures of polarized Tyurin degenerations,
including the geometric identification of the two \(E_8\)-summands of
\(\operatorname{Gr}^W_2(T)\) with the two rational elliptic components
(GT24, §3); the existence of compatible fibrations with allowable loops on
very general \(H\oplus E_8\oplus E_8\)-polarized K3 surfaces (GT24, §4); and
the mirror-pair criterion for very general data (GT24, §5, stated as
Theorem 6.2 with the content with which it is used). The role of the
corrected positivity condition \((A|_{V_i})^2>0\) — it forces *both*
components to contribute a full \(E_8\), hence exactly two
\(\mathrm{II}^*\) fibres on the mirror — is isolated in Remark 4.6, which
also explains why the statement fails without it.